\chapter{The Business Services and the Common Module}
\label{ch:services}

The previous chapters covered the plumbing services. This one maps the
three that do the actual work, plus the shared \code{common} module ---
and, through them, the two ways microservices talk to each other.

\section{course-service: the workhorse}

Courses, students, assignments, enrollments, submissions --- classic
CRUD over Postgres (Chapter~7), plus file attachments stored in MinIO
(Chapter~9). It \emph{publishes} \code{ts.assignment.created} events and
\emph{listens} for \code{ts.user.activated} --- the pattern to notice is
that it reacts to another service's business event without ever calling
that service.

\section{auth-service: identity}

Everything Chapter~5 described: registration with e-mail activation,
password and Google login, TOTP 2FA, JWT issuing, refresh tokens. It
publishes three events: \code{ts.user.registered},
\code{ts.user.activated}, \code{ts.user.admin-provisioned}.

\section{dictionary-service: the small one}

Vocabulary lookups over MongoDB (Chapter~8). Deliberately simple --- it
exists partly to prove the architecture accommodates a document store and
partly as the low-stakes service to try things on first.

\section{notification-service: the pure consumer}

No REST API of consequence, no database. It subscribes to three topics
and turns events into e-mails (Maildev in dev, a real SMTP relay later):

\begin{center}
\small
\begin{tabular}{@{}lll@{}}
\toprule
Topic & Producer & Mail sent \\
\midrule
\code{ts.user.registered} & auth-service & activation link \\
\code{ts.user.admin-provisioned} & auth-service & credentials notice \\
\code{ts.assignment.created} & course-service & new-assignment notice \\
\bottomrule
\end{tabular}
\end{center}

Because its \emph{only} input is Kafka, it is the one service that cannot
usefully run before a broker exists --- which is why its Kubernetes
manifest shipped with the Kafka milestone (Chapter~25) rather than with
the other services.

\section{Synchronous vs.\ asynchronous, in this codebase}

The project uses both interaction styles, each where it belongs:

\begin{description}
  \item[Synchronous (HTTP through the gateway)] when the caller needs the
  answer to proceed: the frontend fetching a course list. Coupled in
  time --- both sides must be up.
  \item[Asynchronous (events through Kafka)] when the caller only needs
  the fact recorded: a registration happened; whoever cares can act.
  auth-service does not know notification-service exists. Decoupled in
  time --- mail goes out late, not never, if the consumer is down.
\end{description}

The rule of thumb the codebase follows: \emph{commands} that need answers
are HTTP; \emph{facts} that others may react to are events, named in the
past tense (\code{UserRegistered}, \code{AssignmentCreated}).

\section{The common module}

\begin{lstlisting}
com/ts/common/
  dto/        CourseDto, StudentDto, UserDto, ...
              UserRegisteredEvent, UserActivatedEvent,
              AssignmentCreatedEvent, AdminProvisionedUserEvent
  exception/  ApiException, ErrorResponse
  security/   JwtConstants     <- header names, claim names
\end{lstlisting}

Three kinds of sharing, with three different risk profiles:

\begin{itemize}
  \item \textbf{Event records} are the \emph{contract} between producer
  and consumer. Sharing the class guarantees both sides agree --- at the
  price of lockstep deployment when the shape changes. (The industrial
  alternative --- schema registry --- appears in Chapter~11.)
  \item \textbf{\code{JwtConstants}} shares header and claim \emph{names}
  between gateway and services. A typo'd header name would fail silently
  --- requests simply arrive anonymous --- so centralizing it is cheap
  insurance.
  \item \textbf{DTOs and exceptions} are convenience sharing: pleasant in
  a monorepo, and the first thing to duplicate-instead-of-share if
  services ever split into separate repositories, because shared
  convenience classes are how ``independent'' services end up released
  together.
\end{itemize}

\begin{decision}[a shared module at all?]
Microservice purists forbid shared code between services. The purist
alternative --- each service defines its own copy of every event --- makes
drift possible precisely where agreement is mandatory. This project
draws the line at: share \emph{contracts} (events, header names), think
twice about sharing anything with behavior. \code{jib.skip=true} in
common's POM marks it as the one module that is a library, not a
deployable.
\end{decision}

\begin{pitfall}[the lockstep trap]
Symptom: after adding a field to \code{UserRegisteredEvent} and deploying
only auth-service, notification-service starts throwing deserialization
errors on every message --- and because it shares a consumer group, it
keeps retrying the same broken message. Cause: the contract changed on
one side only, and the topic now holds events in two shapes. Rules that
prevent it: only \emph{add} optional fields, never rename or remove;
deploy consumers before producers; and configure the consumer's error
handling (Chapter~11) so one poison message cannot stall the topic.
\end{pitfall}

\begin{quiz}
\begin{enumerate}
  \item For each pair, say whether HTTP or an event fits and why:
  (a) frontend needs a student list; (b) a submission was graded and the
  student should be told; (c) course-service must verify a JWT.
  \item Why is notification-service the only service whose Kubernetes
  manifest waits for the Kafka milestone, when auth-service also uses
  Kafka?
  \item Name one thing that belongs in \code{common} and one that should
  never be there, with the reasoning.
  \item What deployment-order rule makes an event schema change safe,
  and what property must the change itself have?
\end{enumerate}
\end{quiz}
